% ---------------------------------------------------------------------
%  CHƯƠNG 4 — CÀI ĐẶT VÀ THỬ NGHIỆM
% ---------------------------------------------------------------------
\chapter{Cài đặt và thử nghiệm}
\label{ch:caidat}

\begin{muctieu}
Chương này mô tả môi trường triển khai, các thành phần chính kèm mã nguồn
minh hoạ, quy trình huấn luyện, giao diện demo và cách kiểm thử hệ thống.
\end{muctieu}

\section{Môi trường triển khai}

Hệ thống được phát triển và huấn luyện trên máy tính cá nhân chạy Windows 11,
Python 3.11, PyTorch bản CPU cùng các thư viện ở Bảng~\ref{tab:congcu}. Mạng
nơ-ron có kích thước nhỏ (đầu vào 4 bước, 2 tầng LSTM 128 nơ-ron) nên huấn
luyện được trên CPU, không cần GPU. Phụ thuộc được khai báo trong
\code{pyproject.toml}, chia nhóm \code{train}, \code{api}, \code{ingestion},
\code{dev}.

Mã nguồn được quản lý bằng \code{git} và đặt tại kho:
\begin{center}
\url{https://github.com/VietManh05/typhoon-track-vietnam}
\end{center}

\textbf{Sử dụng công cụ trí tuệ nhân tạo hỗ trợ:} nhóm dùng trợ lý lập trình
Claude (Anthropic) để hỗ trợ viết một phần mã nguồn và soạn thảo báo cáo LaTeX.
Mọi số liệu trong báo cáo được sinh ra bằng cách chạy mã nguồn của dự án, và
nhóm đã kiểm tra lại toàn bộ nội dung.

\section{Các thành phần chính}

\subsection{Mô hình LSTM dự báo độ dịch chuyển}

Mã nguồn~\ref{lst:lstm} là hàm lan truyền xuôi của \code{LSTMTrackForecaster}
(tệp \code{src/typhoon\_vn/models/lstm.py}). Đầu ra hồi quy được cộng với vị trí
quan trắc cuối (đã đổi về độ thật), hiện thực hoá thiết kế dự báo độ dịch chuyển.

\begin{lstlisting}[language=Python, caption={Lan truyền xuôi của mô hình LSTM (rút gọn, bỏ các đầu ra gió và khí áp).}, label={lst:lstm}]
def forward(self, x, mask=None):
    # x: (batch, time, features)
    context = self.encode(x, mask)  # last hidden state
    H = self.n_horizons
    reg = self.reg_head(context).view(-1, H, 2)
    if self.residual:
        _, lengths = _validate_sequence(x, mask)
        # last observed lat/lon, scaled
        last = _last_valid(x[..., :2], lengths)
        # back to physical degrees
        anchor = (last * self.position_scale
                  + self.position_mean)
        # position + predicted displacement
        reg = anchor.unsqueeze(1) + reg
    cls = self.cls_head(context).view(-1, H, self.n_classes)
    return {"reg": reg, "cls": cls}
\end{lstlisting}

\subsection{k-NN analog và lời giải thích}

Mã nguồn~\ref{lst:knn} cài đặt Giải thuật~\ref{alg:knn}
(tệp \code{src/typhoon\_vn/models/classical.py}): lấy $k$ láng giềng, đánh
trọng số nghịch đảo khoảng cách, cộng độ dịch chuyển trung bình vào vị trí cuối.
Hàm \code{explain} trả về các cơn bão lịch sử khác nhau gần nhất để hiển thị.

\begin{lstlisting}[language=Python, caption={Dự báo của k-NN analog.}, label={lst:knn}]
def predict(self, x_raw):
    """Forecast positions (N, H, 2) in degrees."""
    distance, index = self._neighbours(x_raw)
    # inverse-distance weights, normalised per forecast
    weight = 1.0 / (distance + 1e-6)
    weight /= weight.sum(axis=1, keepdims=True)
    targets = self._targets[index]
    displacement = np.einsum("nk,nkhc->nhc", weight, targets)
    return anchors(x_raw)[:, None, :] + displacement
\end{lstlisting}

\subsection{Vòng huấn luyện}

Lớp \code{Trainer} (tệp \code{src/typhoon\_vn/training/trainer.py}) thực hiện
vòng lặp epoch, cắt gradient, lịch học cosine, dừng sớm theo validation, lưu
checkpoint và cho phép huấn luyện tiếp từ checkpoint (khôi phục cả trạng thái
bộ sinh số ngẫu nhiên để tái lập kết quả).

\subsection{Đóng gói mô hình}

Sau huấn luyện, mô hình được đóng gói thành thư mục gồm \code{weights.pt},
\code{scaler.json}, \code{split-manifest.json} và \code{manifest.json}. Tệp
manifest lưu phiên bản, mã băm SHA-256 của tập dữ liệu, cấu hình mô hình, mã
băm danh sách đặc trưng, bán kính nón và checksum từng tệp. Khi nạp, hệ thống
kiểm tra toàn bộ checksum, số đặc trưng, số mốc dự báo và nạp trọng số với
\code{weights\_only=True}; gói bị sửa hoặc sai định dạng sẽ bị từ chối.

\subsection{Chạy các thí nghiệm}

Mã nguồn~\ref{lst:run} là các lệnh tái lập toàn bộ kết quả ở
Chương~\ref{ch:ketqua}.

\begin{lstlisting}[language=bash, caption={Các lệnh tái lập thí nghiệm.}, label={lst:run}]
# Benchmark on IBTrACS and uncertainty-cone calibration
python scripts/run_ibtracs_benchmark.py --data-root data \
    --output-dir runs/ibtracs-benchmark
# Compare all methods with three seeds
python scripts/run_model_comparison.py --data-root data \
    --output-dir runs/model-comparison --seeds 0 1 2
# Error-by-horizon figure and case studies
python scripts/plot_error_analysis.py --run runs/model-comparison \
    --data-root data
# Build the analog library used by the API
python scripts/build_analog_artifact.py --data-root data \
    --output runs/analog/analog-wp.npz
\end{lstlisting}

\section{Giao diện và demo}

Mô hình được phục vụ qua FastAPI với các điểm cuối chính \code{/forecast},
\code{/typhoons/active} và \code{/typhoons/\{id\}/track} (quỹ đạo dự báo kèm
nón bất định). Giao diện bản đồ web hiển thị quỹ đạo thực, quỹ đạo dự báo của
LSTM, nón bất định, quỹ đạo gợi ý của k-NN analog (đường màu tím) và khung
``Gợi ý theo bão tương tự'' liệt kê 5 cơn bão lịch sử giống nhất. Thư viện
analog chỉ gồm bão 1980--2020, không chứa bão kiểm thử.

Bản demo công khai tại \url{https://baoviet2026.pages.dev} phát lại 20 cơn bão
thật giai đoạn 2021--2024 ảnh hưởng Việt Nam, tại thời điểm vị trí +72 giờ gần
bờ biển nhất, kèm quỹ đạo thực tế và sai số. Hình~\ref{fig:yagi} là ví dụ cho
bão YAGI 2024.

\begin{figure}[H]
\centering
\includegraphics[width=0.85\textwidth]{case_WP122024.png}
\caption{Quỹ đạo dự báo của các phương pháp cho bão YAGI (phát lúc 06Z 04/09/2024).}
\label{fig:yagi}
\end{figure}

\begin{luuy}
Khi phát lại một cơn bão cũ với thời gian dời về hiện tại, đặc trưng mùa của
k-NN thay đổi (ví dụ YAGI tháng 9 thành tháng 11) và cho ra bão tương tự khác
hẳn. Vì vậy bản demo giữ đúng ngày lịch sử của cơn bão.
\end{luuy}

\section{Kiểm thử}

Dự án có bộ kiểm thử tự động bằng \code{pytest} (thư mục \code{tests/}), chia
thành ba nhóm:
\begin{itemize}[leftmargin=1.6em, itemsep=2pt]
  \item \textbf{Dữ liệu:} parser các nguồn, chuẩn hóa mã bão, gộp đa nguồn, ghi
        Parquet; dữ liệu biên như tọa độ ngoài phạm vi, thời điểm thiếu múi giờ,
        giá trị gió không hữu hạn đều bị từ chối.
  \item \textbf{Học máy:} bộ chuẩn hóa chỉ fit trên train, chia tập không trùng
        cơn bão, kích thước đầu ra mô hình, huấn luyện tiếp từ checkpoint, gói mô
        hình bị sửa checksum phải bị từ chối.
  \item \textbf{API:} kiểm tra đầu vào (ít nhất 2 quan trắc, khoảng cách thời
        gian hợp lệ, tốc độ suy ra không quá 150~km/h) và các điểm cuối.
\end{itemize}

Một lỗi phát hiện trong quá trình làm là timestamp trùng hoặc lùi trong dữ
liệu gốc sinh ra tốc độ và gia tốc giả rất lớn; lỗi được khắc phục bằng cách
không gán khoảng thời gian 6 giờ giả định cho các bước này.
